% Related Work -- draft v1 (2026-08-24)
% \cite keys are placeholders; fill in refs.bib. Key -> paper mapping at bottom.

\section{Related Work}
\label{sec:related}

\subsection{Clinical Anchoring and Misleading Context}
\label{sec:related-anchoring}

Diagnostic reasoning rarely begins from an unframed case. Referral letters
can carry a provisional diagnosis from an upstream clinician, and controlled
human studies show that such suggestions can narrow the differential or
anchor subsequent judgments \cite{spaanjaars2015referral, staal2022referral}.
This motivates our intervention as a clinically plausible
\emph{referral-mediated} anchoring scenario. It does not imply that every
referral contains a diagnosis or that every clinical model is deployed in
this workflow.

Medical-LLM robustness work has already established the behavioral problem.
BiasMedQA injects clinically motivated cognitive-bias statements into USMLE
questions \cite{schmidgall2024biasmedqa}. MED-STRESS studies abandonment of
initially correct diagnoses under escalating multi-turn pressure
\cite{xiao2026medstress}, while MedMisBench measures large accuracy drops
under misleading clinical context \cite{zhou2026medmisbench}. Narrative
Anchoring holds clinical facts fixed while varying sociolinguistic register
\cite{singh2026narrative}. These studies show that benchmark knowledge does
not guarantee resilience to context; the average behavioral drop is
therefore not our novelty.

Our question begins after that drop. We use a within-case,
placebo-controlled four-arm intervention in which the patient findings are
fixed and only the referral sentence varies. This separates sentence
insertion from the content-specific cost of a wrong suggestion, yields a
per-case counterfactual label for whether the answer moved, and lets us ask
where gold, suggested, and third-diagnosis signals go inside the model. We
replicate the behavioral effect in case-report language, but claim the
internal mechanism only on the controlled DDXPlus testbed.

\subsection{Chain-of-Thought Faithfulness and Internal--Output Dissociation}
\label{sec:related-dissociation}

Generated reasoning is useful behaviorally but is not automatically a causal
record of how an answer was produced. Turpin et al. show that features that
move answers can go unmentioned while a model rationalizes the biased answer
\cite{turpin2023say}. Lanham et al. intervene on reasoning traces and find
that dependence on CoT varies across tasks and models
\cite{lanham2023measuring}. In medicine, Afolabi et al. use causal ablation,
positional perturbation, and hint injection on closed-source assistants and
find that external suggestions can be incorporated without acknowledgment
\cite{afolabi2026faithful}. These results motivate comparing CoT with an
internal channel; they do not justify assuming that CoT contains no signal.
We therefore evaluate both rule-based features and a strong LLM monitor.

The closest single-run attribution study is \emph{Catching Rationalization}
\cite{mirtaheri2026rationalization}. In general-domain multiple-choice tasks,
pre-generation probes match a full-CoT LLM monitor and post-generation probes
outperform it. Our causal label differs: their hint points to a listed option,
whereas most moved outputs in our open-diagnosis setting go to a third
diagnosis rather than copy the suggestion. We ask whether one wrong-note run
identifies cases that a hidden no-note counterfactual shows were causally
moved, including a subset where output copying is blind by construction.

Medical studies also show that activation content can exceed what a model
states. Fraile Navarro et al. use the same released Gemma-3-12B NLA and
layer-32 activations to localize a triage output-format failure
\cite{frailenavarro2026internal}. Tayebi Arasteh recovers evidence grades from
hidden states when the stated grades are near chance
\cite{tayebi2026evidence}. Basu et al. report strong clinical-risk probe
performance despite substantially lower output sensitivity
\cite{basu2026risk}. These are direct precedents, not gaps we claim to fill.
We differ by studying final diagnosis under a referral-note intervention,
resolving gold/suggestion/other trajectories at prompt landmarks, and
connecting case-level causal attribution to a controlled correction ladder.

\subsection{Reading and Acting on Activations}
\label{sec:related-internals}

Linear probes recover prespecified variables from hidden states but establish
decodability rather than causal use \cite{belinkov2022probing}. Logit and tuned
lenses decode token distributions \cite{belrose2023tunedlens}, while sparse
autoencoders decompose activations into learned feature dictionaries
\cite{cunningham2023sae, bricken2023monosemanticity}. These tools suit closed
questions and are our primary quantitative instrument on DDXPlus.

Natural-language readouts instead aim to express open-ended activation
content. Patchscopes, SelfIE, and LatentQA decode hidden representations into
text \cite{ghandeharioun2024patchscopes, chen2024selfie,
pan2024latentqa}. Natural Language Autoencoders jointly train an activation
verbalizer and reconstructor through a natural-language bottleneck
\cite{anthropic2026nla}. NLAs can surface unverbalized content and support
hypothesis generation, but their authors also report confabulation. More
generally, Li et al. show that activation-verbalization benchmarks can often
be solved without target-model internals and that descriptions may reflect
the verbalizer's parametric knowledge \cite{li2026privileged}. We therefore
do not infer faithfulness from fluent text or SFT loss. Our activation
verbalizer is admitted as a measurement only after matched-versus-shuffled,
swap, heldout-content, memorization, and cross-case-contamination controls.
Even then, we report it as a complementary open-vocabulary channel, not a
replacement for the stronger supervised probe or a validated
clinician-facing explanation.

Finally, decodability does not guarantee controllability. Selective
re-prompting can use internal error signals while limiting disruption to
correct outputs \cite{sun2025selective}. In contrast, recent work reports
gaps between successful detection and activation-level correction
\cite{basu2026risk, vankadaru2026hallucination, liu2026decodability}. We test a
different deployment path: leave source-model weights and activations
untouched, externalize a decoded content candidate, and feed it back as text.
Our result supports only a conditional claim: accurate content can recover
moved answers, but indiscriminate re-prompting damages correct ones, and no
independent advantage of natural-language form remains after content
accuracy is controlled.

% ---------------------------------------------------------------------------
% \cite key -> paper mapping (fill refs.bib from this list; delete when done)
% ---------------------------------------------------------------------------
% wang2024safety            Safety challenges of AI in medicine in the era of LLMs. arXiv:2409.18968
% asgari2025hallucination   Framework to assess clinical safety and hallucination rates (med summarisation). npj Digit Med, s41746-025-01670-7
% ethics2025review          Systematic review of ethical considerations of LLMs in healthcare. PMC12460403
% chen2023drift             How is ChatGPT's behavior changing over time? arXiv:2307.09009  [optional]
% clinicalcode2025          Cracking the clinical code: scoping review on mechanistic interpretability in medical report generation. ScienceDirect S2950363925000377
% clinicalreasoningfails2026 Why LLMs' Clinical Reasoning Fails. medRxiv 2026.01.26.26344845
% frailenavarro2026internal  Internal Representation, Not Clinical Knowledge. arXiv:2605.29889
% tayebi2026evidence         Internal-verbalized evidence-grade dissociation. arXiv:2606.29034
% basu2026risk               Clinical risk probe and knowledge-action gap. arXiv:2603.18353
% spaanjaars2015referral     Referral diagnosis and anchoring. DOI:10.1027/1015-5759/a000235
% staal2022referral          Referral information and diagnostic reasoning. BMC Med Educ 2022
% xiao2026medstress          MED-STRESS. arXiv:2605.23932
% zhou2026medmisbench        MedMisBench. arXiv:2606.12291
% singh2026narrative         Narrative Anchoring. arXiv:2607.27384
% afolabi2026faithful        Faithful or Just Plausible? arXiv:2603.13988
% mirtaheri2026rationalization Catching Rationalization. arXiv:2603.17199
% sun2025selective           Selective reprompting with internal error probes. arXiv:2507.12379
% vankadaru2026hallucination Medical hallucination detection/control gap. arXiv:2607.00158
% liu2026decodability        Decodability-control gap. arXiv:2605.05715
% croskerry2003             The importance of cognitive errors in diagnosis... Acad Med 2003
% mahajan2025bias           Cognitive bias in clinical LLMs. npj Digit Med 8:428 (2025)
% lundberg2017shap          A unified approach to interpreting model predictions (SHAP). NeurIPS 2017, arXiv:1705.07874
% ribeiro2016lime           "Why should I trust you?" (LIME). KDD 2016, arXiv:1602.04938
% xaihealthcare2026         XAI in healthcare use-case review. Frontiers AI 2026 (frai.2026.1749527)
% cbm2026                   Concept bottleneck models in medical imaging (e.g., arXiv:2410.15446 PMLR 2026; or radiologist-guided arXiv:2605.07785)
% turpin2023say             Language models don't always say what they think. NeurIPS 2023, arXiv:2305.04388
% lanham2023measuring       Measuring faithfulness in chain-of-thought reasoning. arXiv:2307.13702
% chen2025reasoning         Reasoning models don't always say what they think. Anthropic 2025, arXiv:2505.05410
% schmidgall2024biasmedqa   BiasMedQA: cognitive biases in medical LLMs. npj Digit Med 2024, arXiv:2402.08113
% reasoningnoprotect2025    LLM reasoning does not protect against clinical cognitive biases. medRxiv 2025 (저자 확인)
% faithfulplausible2025     Faithful or just plausible? Faithfulness of closed-source LLMs in medical reasoning. NeurIPS 2025, arXiv:2603.13988
% clinicalgraphs2026        Clinical reasoning graphs. arXiv:2606.29876
% belinkov2022probing       Probing classifiers: promises, shortcomings, advances. CL 2022
% nostalgebraist2020logitlens  Interpreting GPT: the logit lens. LessWrong 2020
% belrose2023tunedlens      Eliciting latent predictions with the tuned lens. arXiv:2303.08112
% cunningham2023sae         Sparse autoencoders find highly interpretable features. arXiv:2309.08600
% bricken2023monosemanticity Towards monosemanticity. Transformer Circuits 2023
% ghandeharioun2024patchscopes Patchscopes. ICML 2024, arXiv:2401.06102
% chen2024selfie            SelfIE. ICML 2024, arXiv:2403.10949
% pan2024latentqa           LatentQA. arXiv:2412.08686
% anthropic2026nla          Natural language autoencoders produce unsupervised explanations of LLM activations. Transformer Circuits 2026
% li2026privileged          Do activation verbalization methods convey privileged information? ICML 2026, arXiv:2509.13316
% rationalization2026       Catching rationalization in the act: activation probing. arXiv:2603.17199
% truthoverridden2026       When truth is overridden: internal origins of sycophancy. AAAI 2026, arXiv:2508.02087
% yuan2026hidden            Hidden error awareness in CoT: diagnostic, not causal. arXiv:2605.09502
% mehrafarin2026hidden      When CoT fails, the solution hides in the hidden states. arXiv:2604.23351
% adrprobing2025            Probing LLM hidden states for adverse drug reaction knowledge. PMC11844579
% alignmentresistant2025    Probing hidden states for calibrated, alignment-resistant predictions. medRxiv 2025.09.17.25336018
% jmirsae2026               SAEs for mechanistic interpretability of LLMs in medicine. JMIR AI 2026, e81134
% ehrsae2026                SAE decomposition of clinical sequence model representations. arXiv:2605.04072
